% Generated by Listas/make_statements.py from pratica-lecture-03.md, item 12. Do not edit.
\textbf{Source.} Tadelis (2013), \emph{Game Theory: An Introduction}, Exercise 8.10 ``Playing It Safe'', p.~172 {[}PDF 191{]}. Reproduced verbatim. Answer in the source: none in the book (the preface, p.~xiii, says about half are solved on Princeton UP's website; not checked).



\begin{quote}

\textbf{Playing It Safe:} Consider the following dynamic game: Player 1 can choose to play it safe (denote this choice by \(S\)), in which case both he and player 2 get a payoff of 3 each, or he can risk playing a game with player 2 (denote this choice by \(R\)). If he chooses \(R\) then they play the following simultaneous-move game:



{\def\LTcaptype{none} % do not increment counter

\begin{longtable}[]{@{}ccc@{}}

\toprule\noalign{}

Player 1 ~Player 2 & \(A\) & \(B\) \\

\midrule\noalign{}

\endhead

\bottomrule\noalign{}

\endlastfoot

\(C\) & 8, 0 & 0, 2 \\

\(D\) & 6, 6 & 2, 2 \\

\end{longtable}

}



\begin{enumerate}

\def\labelenumi{\alph{enumi}.}

\item

  Draw a game tree that represents this game. How many proper subgames does it have?

\item

  Are there other game trees that would work? Explain briefly.

\item

  Construct the matrix representation of the normal form of this dynamic game.

\item

  Find all the Nash and subgame-perfect equilibria of the dynamic game.

\end{enumerate}

\end{quote}

